\documentclass{article}

\input{header}

\title{CSCI 2590: Assignment 4}

\author{Keefer Wu \\ N19122827 \\ \texttt{https://github.com/KeeferW}}

\date{}


\colmfinalcopy
\begin{document}
\maketitle

\section*{Q0. 1.}
GitHub repository link: \texttt{https://github.com/KeeferW/NLP-HW4}

\section*{Q2. 1.}

The transformation does synonym replacement using WordNet. Each word in the review gets replaced with a synonym at probability 0.15. The text is tokenized with NLTK's \texttt{word\_tokenize}, and for each token we sample a Bernoulli trial. If selected, we look up synsets for that word, collect lemmas from the first synset that are different from the original word, and pick one at random. Words with no synsets or no alternate lemmas stay unchanged. The tokens are then rejoined with \texttt{TreebankWordDetokenizer}.

This is a reasonable transformation because swapping words for synonyms does not change the sentiment of a review. A human would give the same label to ``The film was excellent'' and ``The film was superb.'' It also reflects real test time variation, since users naturally use different vocabulary then what appears in training data.

\section*{Q3. 1}
\textbf{Report \& Analysis}
    \begin{itemize}
        \item The BERT model trained on the original data scored 90.90\% on the original test set and 83.82\% on the transformed test set. After training with augmented data, it scored 92.93\% on the original test set and 86.59\% on the transformed test set.

        \item Performance on the transformed set improved after augmentation from 83.82\% to 86.59\%, a gain of 2.77 points. The drop from original to transformed accuracy was 7.08\% for the base model versus 6.34\% for the augmented model, showing the augmented model is more robust to the transformation.

        \item The improvement on the transformed set suggests that exposing the model to synonym-replaced and typo-perturbed examples during training reduces its sensitivity to those perturbations at test time. Since the transformations preserve label meaning, the augmented examples add useful signal without introducing noise.

        \item The main limitation is that this approach only helps with the specific transformations it was trained on. If test inputs differ in some other way such as a different domain, writing style, or longer texts, the augmented model offers no advantage over the baseline.
    \end{itemize}

\section*{Part II. Q4}

\begin{table}[h!]
\centering
\begin{tabular}{lcc}
\toprule
Statistics Name & Train & Dev \\
\midrule
Number of examples & 4225 & 466 \\
Mean sentence length (tokens) & 18.1 & 18.1 \\
Mean SQL query length (tokens) & 217.4 & 211.1 \\
Vocabulary size (natural language) & 792 & 466 \\
Vocabulary size (SQL) & 556 & 396 \\
\bottomrule
\end{tabular}
\caption{Data statistics before any pre-processing. Token counts use the T5 subword tokenizer.}
\label{tab:data_stats_before}
\end{table}

\begin{table}[h!]
\centering
\begin{tabular}{lcc}
\toprule
\textbf{T5 fine-tuned model} & Train & Dev \\
\midrule
Mean encoder input length (tokens) & 23.1 & 23.1 \\
Mean decoder target length (tokens) & 218.4 & 212.1 \\
Vocabulary size (encoder inputs) & 811 & 483 \\
Vocabulary size (decoder targets) & 557 & 397 \\
\midrule
\bottomrule
\end{tabular}
\caption{Data statistics after pre-processing. The task prefix adds a few tokens to each encoder input.}
\label{tab:data_stats_after}
\end{table}

\newpage

\section*{Q5}\label{sec:t5}

\begin{table}[h!]
\centering
\begin{tabular}{p{3.5cm}p{10cm}}
\toprule
Design choice & Description \\
\midrule
Data processing & A compact database schema is prepended to each NL query in the format \texttt{"translate English to SQL: tables: flight(...), airport(...), ... | <query>"}. This gives the model explicit knowledge of available tables and columns. \\
Tokenization & Used the default \texttt{google-t5/t5-small} tokenizer (\texttt{T5TokenizerFast}) for both encoder and decoder. Encoder input has EOS appended. Decoder input starts with the pad token (id=0, which T5 uses as the decoder start token) followed by the SQL tokens. Decoder target is the SQL tokens followed by EOS. Padding is done dynamically each batch. \\
Architecture & The full pretrained \texttt{google-t5/t5-small} encoder-decoder was fine-tuned end to end with no layers frozen. Inference uses beam search with 4 beams, \texttt{max\_new\_tokens=350}, and \texttt{early\_stopping=True}. \\
Hyperparameters & Learning rate: $5 \times 10^{-4}$, AdamW optimizer ($\beta_1$=0.9, $\beta_2$=0.999), weight decay: 0.01, cosine scheduler with 1 warmup epoch, batch size: 16, max 25 epochs, early stopping with patience 5 based on dev Record F1. \\
\bottomrule
\end{tabular}
\caption{Details of the best-performing T5 model configurations (fine-tuned)}
\label{tab:t5_results_ft}
\end{table}

\section*{Q6.}
\label{Q6}

\paragraph{Quantitative Results:}
\begin{table}[h!]
\centering
\begin{tabular}{lcc}
  \toprule
  System & Query EM & F1 score \\
  \midrule
  \multicolumn{3}{l}{\textbf{Dev Results}} \\
  \midrule
  \multicolumn{3}{l}{T5 fine-tuned} \\
  Full model & 3.22 & 76.93 \\
  \midrule
  \multicolumn{3}{l}{\textbf{Test Results}} \\
  \midrule
  T5 fine-tuning & -- & -- \\
  \bottomrule
\end{tabular}
\caption{Development and test results. Values are percentages.}
\label{tab:results}
\end{table}

\paragraph{Qualitative Error Analysis:}

\begin{table}[h!]
  \centering
  \small
  \setlength{\tabcolsep}{4pt}
  \begin{tabular}{p{2.0cm}p{3.2cm}p{4.0cm}p{2.8cm}}
    \toprule
    \textbf{Error Type} & \textbf{Example} & \textbf{Description} & \textbf{Statistics} \\
    \midrule
    Wrong table join &
      NL: ``flights from Boston on Monday''; query omits \texttt{days} join &
      Valid SQL but missing joins for temporal constraints; returns incorrect result set &
      $\sim$25/466 queries \\
    \midrule
    Incorrect city mapping &
      NL: ``cheapest fare to Dallas''; \texttt{city\_name} filter absent &
      City names not translated into \texttt{WHERE} clause; geographic constraint ignored &
      $\sim$18/466 queries \\
    \midrule
    SQL syntax error &
      NL: complex multi-hop query; model generates malformed SQL &
      14.16\% of outputs are syntactically invalid; returns empty record set &
      $\sim$66/466 queries \\
    \bottomrule
  \end{tabular}
  \caption{Qualitative error analysis on the dev set for the fine-tuned T5 model.}
  \label{tab:qualitative}
\end{table}

\section*{Q7.}

Google Drive link to model checkpoint:
 https://drive.google.com/drive/folders/1VA0VNok0G-zsLas9mfHTShEhxeihnKVa?usp=sharing

\section*{Extra Credit 2-1: T5 from Scratch}

Same setup as fine-tuning but with randomly initialized weights. Used a slightly lower learning rate of $1 \times 10^{-3}$ since the model is less stable at the start of training. Performance is lower then fine-tuning as expected since there are no pretrained weights to build on.

\begin{table}[h!]
\centering
\begin{tabular}{lcc}
  \toprule
  System & Query EM & F1 score \\
  \midrule
  \multicolumn{3}{l}{\textbf{Dev Results}} \\
  \midrule
  T5 from scratch & XX.XX & XX.XX \\
  \midrule
  \multicolumn{3}{l}{\textbf{Test Results}} \\
  \midrule
  T5 from scratch & XX.XX & XX.XX \\
  \bottomrule
\end{tabular}
\caption{Results for T5 trained from scratch.}
\end{table}

Google Drive link to from-scratch checkpoint: \texttt{<TODO>}

\section*{AI Usage}

\noindent\textbf{1. AI Tools Used} \\
ChatGPT \\
\hrulefill

\noindent\textbf{2. How AI Was Used} \\
Used to get clarification on technical concepts and to assist with writing the report. \\
\hrulefill

\noindent\textbf{3. Example Prompts / Interactions} \\
(1) ``Explain the architecture of T5 and how it can be fine-tuned for sequence-to-sequence tasks like text-to-SQL.'' \\
(2) ``Given the following accuracy values help me fill in the results table in LaTeX format'' \\
(3) ``In simple chronological steps, how would I approach implementing an example T5Dataset with encoder/decoder tokenization using T5TokenizerFast?'' \\
\hrulefill

\end{document}